\section{Well-posedness of the heat equation}
\label{sec:wellposedness}

\subsection*{Abstract standing assumption}

Let $T>0$ be the (fixed) control time horizon, and let $H$ and $U$ be real
Hilbert spaces -- the \emph{state space} and \emph{control space},
respectively (specialized to concrete function spaces below, e.g.\
$H=U=L^2(0,L)$ for internal control, where $L>0$ is the length of the spatial
interval). Controllability problems are then posed abstractly as
\begin{equation}
  \label{eq:abstract-system}
  y'(t) + Ay(t) = Bu(t) \ \text{in } (0,T), \qquad y(0) = y_0 \in H,
\end{equation}
where $u \in L^2(0,T;U)$ is the control, and $A, B$ are (possibly unbounded)
linear operators densely defined on $D(A) \subset H$ and $D(B) \subset U$. It
is assumed throughout that \eqref{eq:abstract-system} is \emph{well posed}: for
every $(y_0, u) \in H \times L^2(0,T;U)$ there is a unique solution
$y \in C([0,T];H)$ depending continuously on the data, understood in the weak
sense via the semigroup generated by $-A$ \citep{pazy2012semigroups}.

\subsection{Internal (distributed) control}
\label{ssec:wp-internal}

For internal control, $H = U = L^2(0,L)$, $A = -\alpha\partial_{xx}$ for a
fixed diffusion coefficient $\alpha>0$, with
$D(A) = H^2(0,L)\cap H^1_0(0,L)$, and $B$ is the (bounded) multiplication
operator by $\mathbf{1}_\omega$ for a subset $\omega \Subset (0,L)$:
\begin{equation}
  \label{eq:heat-internal}
  \begin{cases}
    y_t - \alpha y_{xx} = v(x,t)\mathbf{1}_\omega(x), & (x,t) \in (0,L)\times(0,T),\\
    y(0,t) = y(L,t) = 0, & t \in (0,T),\\
    y(x,0) = y_0(x), & x \in (0,L).
  \end{cases}
\end{equation}

\begin{theorem}[Well-posedness, internal control]
  \label{thm:wp-internal}
  For every $y_0 \in L^2(0,L)$ and $v \in L^2\big((0,L)\times(0,T)\big)$,
  problem \eqref{eq:heat-internal} has a unique solution
  $y \in C\big([0,T];L^2(0,L)\big)$. Moreover, \eqref{eq:heat-internal} is null
  controllable \citep{fursikov1996imanuvilov}.
\end{theorem}

\subsection{Boundary control}
\label{ssec:wp-boundary}

For boundary control, the same equation is instead driven through a Dirichlet
datum at one endpoint:
\begin{equation}
  \label{eq:heat-boundary}
  \begin{cases}
    y_t - \alpha y_{xx} = 0, & (x,t) \in (0,L)\times(0,T),\\
    y(0,t) = 0, \quad y(L,t) = v(t), & t \in (0,T),\\
    y(x,0) = y_0(x), & x \in (0,L).
  \end{cases}
\end{equation}
Because the control enters through the boundary trace rather than as a bounded
source term, \eqref{eq:heat-boundary} only admits a solution \emph{by
transposition}, and in a weaker space than the internal case:

\begin{theorem}[Well-posedness, boundary control]
  \label{thm:wp-boundary}
  For every $y_0 \in L^2(0,L)$ and $v \in L^2(0,T)$, problem
  \eqref{eq:heat-boundary} has a unique weak (transposition) solution
  $y \in C\big([0,T]; H^{-1}(0,L)\big)$, and \eqref{eq:heat-boundary} is null
  controllable \citep{fursikov1996imanuvilov}.
\end{theorem}

From here on, $\Omega$ denotes the spatial domain of a result stated at the
general (possibly multi-dimensional) level of the underlying PDE-control
literature; in the concrete one-dimensional setting of
\eqref{eq:heat-internal}--\eqref{eq:heat-boundary} used throughout the rest of
this document, $\Omega = (0,L)$. Correspondingly, $Q := \Omega\times(0,T)$ is
the space-time cylinder and $\Sigma := \partial\Omega\times(0,T)$ is its
lateral boundary (where the boundary control $v(t)$ of
\eqref{eq:heat-boundary} acts).

\begin{remark}
  The objectives of exact/null controllability require $y(\cdot,T) \in
  L^2(\Omega)$, a strictly stronger requirement than the ambient
  well-posedness class $y \in C([0,T];H^{-1}(\Omega))$ \citep{belgacem2011dirichlet}.
  This gap between the natural well-posedness space and the space in which
  controllability is sought is a recurring theme (see the ill-posedness
  discussion in Section~\ref{sec:known-issues}).
\end{remark}

\subsubsection*{Solution-by-transposition construction}

The rigorous justification for Theorem~\ref{thm:wp-boundary} splits the
solution as $y = z + w$, where $w$ solves an auxiliary problem with the same
data but homogeneous boundary conditions, and $z$ is defined by transposition
against a smooth adjoint problem.

\begin{proposition}
  \label{prop:transposition}
  The auxiliary solution $w$ satisfies
  $w \in L^2(0,T;L^2(\Omega)) \cap C^0([0,T];H^{-1}(\Omega))$
  \citep{evans2010partial}. The adjoint problem
  \[
    -p_t - \Delta p = f \ \text{in } Q, \qquad p = 0 \ \text{on } \Sigma,
    \qquad p(x,T) = 0 \ \text{in } \Omega,
  \]
  with $f \in L^2(0,T;L^2(\Omega))$, has a unique solution
  $p \in L^2\big(0,T; H^1_0(\Omega)\cap H^2(\Omega)\big)$ with
  $p_t \in L^2(Q)$; the transposition identity built from $p$ then yields
  $z \in C^0([0,T];H^{-1}(\Omega))$, so $y = z+w$ is well defined
  \citep{lions2012non}.
\end{proposition}

This construction is also the continuous-theory root of any numerical HUM
variant for boundary control that works directly with the (rougher)
transposition solution rather than with the adjoint state.

\begin{remark}[Gap in the thesis, now filled independently -- not from the
  thesis]
  \label{rem:energy-wellposedness}
  No statement of well-posedness in the classical parabolic energy form
  $y \in C([0,T];L^2(\Omega)) \cap L^2(0,T;H^1_0(\Omega))$ appears in the
  thesis; it consistently works either in the abstract semigroup framework
  above, or in the transposition/weak-solution framework in $H^{-1}$. For
  completeness, the standard energy estimate is:
  \begin{theorem}[Classical energy estimate]
    For $y_0 \in L^2(\Omega)$ and $f \in L^2(0,T;L^2(\Omega))$, the linear heat
    equation $y_t - \Delta y = f$ in $Q$, $y=0$ on $\Sigma$, $y(\cdot,0)=y_0$,
    has a unique weak solution
    $y \in C([0,T];L^2(\Omega)) \cap L^2(0,T;H^1_0(\Omega))$, with
    $y_t \in L^2(0,T;H^{-1}(\Omega))$, obtained via the Galerkin method and the
    energy estimate
    \[
      \sup_{t\in[0,T]}\|y(t)\|_{L^2(\Omega)}^2
      + \int_0^T \|y(t)\|_{H^1_0(\Omega)}^2\,dt
      \;\le\; C\Big(\|y_0\|_{L^2(\Omega)}^2 + \|f\|_{L^2(Q)}^2\Big).
    \]
  \end{theorem}
  \citep[\S 7.1]{evans2010partial}. This is the same book already cited in
  Proposition~\ref{prop:transposition} for a different, thesis-attributed
  result; unlike every thesis-sourced statement in this document (checked
  line-by-line against the primary source), this statement was not verified
  page-by-page against Evans' book -- it is standard textbook material, not a
  literal quote.
\end{remark}
